../cictikz/src/cictikz/data/tex/cictikz_preamble.tex

% Two ways of stopping the on resistance from depending on the signal.
%
% Left: switch the bulks. Both wells follow the input while the gate is on, so
% neither device sees a body effect, and both are returned to their supply rail
% when the gate is off so the junctions stay reverse biased. It is cheap, and
% sometimes it is enough.
%
% Right: bootstrap the gate. The capacitor is charged to V_DD while the switch
% is off and then floated on top of the input, so V_GS is the supply no matter
% where the input sits. The block itself is in tikz/boot_lib.tex, because
% l5_sw3 draws two more of them.
%
% The original leaves the far ends of the two bulk networks as bare symbols.
% They are the supply rails: the n-well goes to V_DD and the p-substrate to
% ground when the switch is off, which is the only bias that keeps both
% junctions reverse biased.

\begin{circuitikz}[american, transform shape]

  ../cictikz/src/cictikz/data/tex/boot_lib.tex

  % ================= Left: transmission gate with switched bulks =========

  % ---- The two pass devices ----
  \draw[thick] (-0.9,1.6) to [Tpmos, n=TP, bulk] (0.9,1.6);
  \draw[thick] (TP.gate) -- (0,3.0);
  \bootPort{0,3.1}{south}{$\overline{c}$}
  \draw[thick] (-0.9,-1.6) to [Tnmos, n=TN, mirror, bulk] (0.9,-1.6);
  \draw[thick] (TN.gate) -- (0,-3.0);
  \bootPort{0,-3.1}{north}{$c$}

  % ---- Source and drain rails ----
  \draw[thick] (-1.8,1.6) -- (-0.9,1.6);
  \draw[thick] (-1.8,-1.6) -- (-0.9,-1.6);
  \draw[thick] (-1.8,1.6) -- (-1.8,-1.6);
  \draw[thick] (0.9,1.6) -- (1.8,1.6);
  \draw[thick] (0.9,-1.6) -- (1.8,-1.6);
  \draw[thick] (1.8,1.6) -- (1.8,-1.6);
  \draw[thick] (-2.5,0) -- (-1.8,0);
  \fill (-1.8,0) circle (0.07);
  \bootPort{-2.6,0}{east}{$A$\;}
  \draw[thick] (1.8,0) -- (2.5,0);
  \fill (1.8,0) circle (0.07);
  \bootPort{2.6,0}{west}{\;$B$}

  % ---- n-well: to A while the gate is on, to V_DD while it is off ----
  \draw[thick] (TP.bulk) -- (0,0.8);
  \draw[thick] (-1.8,0.8) -- (-1.5,0.8);
  \fill (-1.8,0.8) circle (0.07);
  \bootSwH{-1.5}{0.8}{$c$}
  \draw[thick] (-0.5,0.8) -- (0.5,0.8);
  \fill (0,0.8) circle (0.07);
  \bootSwH{0.5}{0.8}{$\overline{c}$}
  \draw[thick, ->] (1.5,0.8) -- (2.1,0.8) node[anchor=west] {$V_{DD}$};

  % ---- p-substrate: to A while the gate is on, to ground while it is off ----
  \draw[thick] (TN.bulk) -- (0,-0.8);
  \draw[thick] (-1.8,-0.8) -- (-1.5,-0.8);
  \fill (-1.8,-0.8) circle (0.07);
  \bootSwH{-1.5}{-0.8}{$c$}
  \draw[thick] (-0.5,-0.8) -- (0.5,-0.8);
  \fill (0,-0.8) circle (0.07);
  \bootSwH{0.5}{-0.8}{$\overline{c}$}
  \draw[thick] (1.5,-0.8) -- (1.9,-0.8);
  \draw[thick] (1.9,-0.8) node[ground, rotate=90] {};

  % ================= Right: bootstrapped gate =================

  \begin{scope}[shift={(7.2,0)}]
    \bootBlock{MB}
    \bootPort{-3.3,0}{east}{$A$\;}
    \bootPort{6.9,0}{west}{\;$B$}
  \end{scope}

\end{circuitikz}
\end{document}
